%% !TEX root = manuscript.tex

\section*{Remerciements}

Je tiens tout d'abord à remercier ma tutrice en entreprise, Sabine Cruanas, pour son
accompagnement constant tout au long de cette année, sa disponibilité et le partage de ses
connaissances. Son expertise et sa confiance ont été déterminantes dans ma montée en
compétences, aussi bien sur le plan technique que méthodologique.

J'adresse également mes remerciements à ma tutrice universitaire, Mme Trottier, pour son suivi
et ses conseils, ainsi qu'à l'ensemble de l'équipe pédagogique du Master MIASHS. Les méthodes
acquises au fil de ces deux années constituent le socle sur lequel j'ai pu m'appuyer pour mener
mes missions.

Je remercie aussi chaleureusement l'ensemble des équipes d'Angelotti avec lesquelles j'ai
travaillé, en particulier la direction financière, le service juridique et les équipes
commerciales. Leur disponibilité et leurs retours ont été essentiels pour adapter mes
réalisations à leurs besoins réels. C'est en écoutant leurs objections que la plupart de mes
livrables ont pris leur forme définitive.

Je remercie enfin le groupe Angelotti de m'avoir offert cette opportunité d'alternance dans un
environnement professionnel formateur. À toutes et à tous, j'adresse mes sincères remerciements
pour leur confiance et leur soutien.

\clearpage
\section*{Résumé}

Ce mémoire retrace ma deuxième année d'alternance de Master 2 MIASHS au sein de la Holding SPA
du groupe Angelotti, promoteur-aménageur implanté en Occitanie et en région PACA. L'année a été
dominée par un chantier structurant pour l'entreprise : le remplacement de l'ERP historique
Grimmo par la solution SPO, éditée par Salvia.

Deux missions principales structurent ce travail. La première est la reprise des données vers
le nouvel ERP, menée en deux temps : les opérations et leurs budgets au premier semestre, sur
un périmètre de 283 opérations immobilières, puis les tiers, c'est-à-dire les clients
acquéreurs, au second semestre. J'y expose les difficultés rencontrées, notamment la gestion
des identifiants dynamiques de l'interface de programmation et le choix des règles permettant
de décider que deux lignes désignent la même personne. La seconde mission est un projet de Data
Science mené de bout en bout, le « Copilote Financier », qui outille la direction financière
sur trois questions de gestion : le risque de dérive des marges, le rythme d'écoulement des
programmes et le pilotage des prix du stock. S'y ajoutent des activités complémentaires, avec
le maintien du système décisionnel et la contribution au déploiement du SSO, puis un retour
réflexif sur mon organisation, mes compétences, l'adéquation avec ma formation et le
positionnement RSE de l'entreprise.

Au fil de ces missions, une même question a servi de fil conducteur : comment reconstituer et
fiabiliser les données d'une entreprise pour qu'elles deviennent exploitables par ses équipes
métiers, et ce que cette structure de départ autorise ensuite en matière de modélisation.

Au-delà des livrables, cette année m'a permis de mesurer combien le travail de structuration
des données conditionne ce qu'une analyse peut en tirer, et combien la rigueur, la preuve et la
communication comptent autant que la technique.
